%======================================================================
\section{Part II --- Python/kwant notebook (\code{tjunctions\_python\_ed.ipynb})}
%======================================================================

The same geometry, texture and protocols, with exact diagonalization and an exact evolution of the full BdG basis. The gate analysis
follows the construction of the user's CZ, QFT and period-finding notebooks: \code{individual\_MZMs}, \code{fix\_unitary}, logical-state
\code{masks} and the Gauss--Legendre \code{step}.

\subsection{Setup (cells 1--2)}

\code{OMP\_NUM\_THREADS}, \code{OPENBLAS\_NUM\_THREADS} and \code{MKL\_NUM\_THREADS} are set to 1 before NumPy is imported (the comment: for
tiny matrices threaded BLAS only adds contention; it also makes the timings single-thread). \code{DATA\_DIR = "tjunctions\_python\_ed\_data"}
receives figures and \code{.npz} files. Cell 2 defines the Pauli matrices \code{tau\_0, tau\_x, tau\_y, tau\_z} and the Kronecker products
\code{txs0} $=\tau_x\sigma_0$, \code{t0sx}, \code{t0sy}, \code{t0sz}, \code{tzsx}, \code{tzsy}, \code{tzs0}, \code{tysy}, and \code{tysx}, \code{tys0} for the unused spinor branches of \code{PHconj} (first factor Nambu,
second spin).

%----------------------------------------------------------------------
\subsection{Section 1: geometry}

Constants: \code{NLEG = 16} (sites per arm, Julia \code{TArray} with $L = 4$), \code{NCELL = 3}, \code{SPACING = 32},
\code{ARM\_ANG}, \code{ARM\_ORIG}, \code{ARM\_DIR} = $((1,0), (-1,0), (0,1))$.

\begin{function}{center\_x(c), array\_pos(c, s, n), label}
As in Julia: $x_c = (c - 2)\cdot32$ and $\mathbf r(c,s,n) = (x_c,0) + \mathbf o_s + n\hat e_s$ (integers). \code{label} is the dictionary
$(x,y)\mapsto(c,s,n)$ over all 144 sites, used to evaluate textures at kwant sites.
\end{function}

\begin{function}{build\_system\_3T()}
\begin{Verbatim}
def build_system_3T():
    def chain(site, mu, Delta, J, S):
        Sx, Sy, Sz = S[site.tag]
        JS = J * (Sx * t0sx + Sy * t0sy + Sz * t0sz)
        return -mu * tzs0 + Delta * txs0 + JS
    def hop_x(sitei, sitej, t, alpha):
        return -t * tzs0 + 1j * alpha * tzsy
    def hop_y(sitei, sitej, t, alpha):
        return -t * tzs0 - 1j * alpha * tzsx
\end{Verbatim}
Finalized kwant system on a square lattice with four orbitals per site, basis $(c_\uparrow, c_\downarrow, c^\dagger_\downarrow, -c^\dagger_\uparrow)$.
Every site in \code{label} gets the on-site function \code{chain}: $-\mu\tau_z + \Delta\tau_x + J\,\tau_0\,\mathbf S\cdot\boldsymbol\sigma$, with the
moment looked up in the parameter \code{S} (a dict keyed by site tag). \code{HoppingKind} adds $x$ and $y$ bonds between every pair of
existing neighbouring sites, so the junction and inter-cell bonds appear automatically. Output: 144 sites, $x = -48\dots47$, legs at
$x = -32, 0, 32$.
\end{function}

%----------------------------------------------------------------------
\subsection{Section 2: spin texture and protocols}

Line-by-line ports of the Julia code (Part I, section 2): \code{smoothstep}, \code{f0}, classes \code{Exchange12}, \code{Exchange23},
\code{Fusion} with a \code{duration()} method, \code{f\_horizontal\_12}, \code{f\_vertical\_12}, \code{f\_braid\_12}, \code{f\_array\_12}
(Julia \code{f\_array\_braid}), \code{f\_h23}, \code{f\_v23}, \code{f\_array\_23} (which contains the cell-2 map of Julia's \code{f\_cell2\_23}),
and \code{f\_fuse(q, t)} (uses \code{f0} from the global scope and \code{NLEG} instead of the hard-coded 16). The leg maps are identical,
including the clamping of the outer sites ($\max(L_1-n, 0)$, $\min(n, L_2)$) and the linear map of cell 2's left arm.

\begin{function}{spins(f)}
Returns the dict \code{S[(x, y)]} $=(\cos\phi\sin\theta, \sin\phi\sin\theta, \cos\theta)$ with $\theta = (-1)^{x+y}(\pi/2)f(c,s,n)$ and
$\phi = \theta_{\rm arm}+\pi/2$ for all sites; the kwant parameter \code{S}.
\end{function}

%----------------------------------------------------------------------
\subsection{Section 3: Hamiltonian}

\code{params = dict(t=1.0, mu=-7.0, alpha=0.9, Delta=2.4, J=5.8)}.

\begin{function}{ham\_kwant(f, **kw)}
Dense $4N\times4N$ matrix \code{sys.hamiltonian\_submatrix} for texture $f$, with \code{params} overridden by \code{kw}. Used for the
reference checks; \code{H\_noJ = ham\_kwant(f0, J=0.0)} is the texture-independent part.
\end{function}

\begin{function}{exchange\_matrix(f, J), ham(f)}
\code{exchange\_matrix} fills the $4\times4$ diagonal blocks $J(S_x\tau_0\sigma_x + S_y\tau_0\sigma_y + S_z\tau_0\sigma_z)$ in kwant's site order
(\code{tags}). \code{ham(f) = H\_noJ + exchange\_matrix(f)}, i.e.\ $H(t) = H_{J=0} + J\sum\tau_0\,\mathbf S\cdot\boldsymbol\sigma$. The markdown
states that this is identical to \code{ham\_kwant} and much faster; the check prints $\|H_{\rm fast} - H_{\rm kwant}\| = 0$, Hermiticity
error 0 and particle--hole error $\|CH^*C + H\| = 0$ with $C = \mathbb1\otimes\tau_y\sigma_y$.
\end{function}

\begin{function}{Cop(v)}
Particle--hole conjugation in the $(u_\uparrow, u_\downarrow, v_\downarrow, -v_\uparrow)$ basis: $C v = (\mathbb1_N\otimes\tau_y\sigma_y)\,v^*$, with
$CH^*C = -H$. Used to check that the columns of \code{U\_fixed} are particle--hole partners.
\end{function}

%----------------------------------------------------------------------
\subsection{Section 4: check against the Julia/TensorBinding array}

Reads \code{tb\_reference/coords.txt} (raises \code{FileNotFoundError} with instructions if the TN notebook has not been run) and builds
a $4N\times4N$ matrix \code{P} that maps kwant's site-major vector to Julia's nambu $\otimes$ spin $\otimes$ site order: for Julia site $j$ at
kwant index $k$, \code{P[b*N + j, 4*k + a] = Wsite[b, a]} with
\[
W_{\rm site} = \begin{pmatrix}1&0&0&0\\0&1&0&0\\0&0&0&-1\\0&0&1&0\end{pmatrix},\qquad
(u_\uparrow, u_\downarrow, v_\uparrow, v_\downarrow)_{\rm Julia} = (a, b, -d, c).
\]

\begin{function}{to\_tb\_basis(H)}
$P H P^T$: the kwant matrix in Julia's basis. Each \code{.bin} file is read as \code{complex128}, reshaped in Fortran order. Results:
\code{initial} $3.5\times10^{-12}$, \code{z\_t1020} $4.3\times10^{-10}$, \code{x100\_t485} $2.0\times10^{-9}$, \code{fuse\_t300}
$1.9\times10^{-11}$, the same numbers as the MPO-vs-dense check of the TN notebook.
\end{function}

%----------------------------------------------------------------------
\subsection{Section 5: subgap LDOS}

Top-level code: \code{la.eigh(H0)} for all eigenpairs; \code{w\_site} sums $|U|^2$ over the four orbitals of each site (so it includes
particle and hole); \code{g\_eta} is the Gaussian with $\eta = 0.01$; $A = w_{\rm site}\,g_\eta$ on 401 points of $[-0.2, 0.2]$. All eigenstates
are included; the markdown notes that those at $|E|\ge0.268$ contribute less than $e^{-23}$ inside the window
($(0.268-0.2)^2/2\eta^2\approx23$). Lowest $|E|$ (one per $\pm$ pair): $0, 0, 0.127, 0.2683$; $\max_{\mathbf r}A(\mathbf r,0) = 15.21$. Saved to
\code{array\_ldos\_kwant.npz} (\code{x, y, omega, A}).

\begin{function}{ldos\_figure(Amat, name, out, vmax)}
Plots $A(\mathbf r, \omega=0)$ on the lattice grid (NaN outside the sites), with the in-plane moments as \code{quiver} arrows and the
$\hat z$ moments as dots; saves to \code{out}. The Python counterpart of the Julia LDOS map (Fig.~\ref{fig:ldos}).
\end{function}

%----------------------------------------------------------------------
\subsection{Section 6: Majoranas, fermion modes and logical states}

\subsubsection*{Copied helper functions}

\begin{function}{dag(A, axis1=-2, axis2=-1)}
Hermitian conjugate: swaps the last two axes and conjugates; for a 1D array only conjugates.
\end{function}

\begin{function}{PHconj(psi, ordering="ltslts", spinor="udd-u")}
\begin{Verbatim}
def PHconj(psi, ordering="ltslts", spinor="udd-u"):
    """Return the particle-hole conjugate C psi = (tau_y sigma_y) psi^* for the kwant spinor "udd-u".
    ..."""
\end{Verbatim}
Particle--hole conjugate of the columns of \code{psi} ($4N_s\times n$). For the site-major ordering \code{"ltslts"} and the kwant spinor
\code{"udd-u"} it applies $\tau_y\sigma_y$ to each site's four components of $\psi^*$ (the case used here). Other branches: spinor \code{"uddu"}
($\tau_y\sigma_x$, \code{tysx}), \code{"udud"} ($\tau_y\sigma_0$, \code{tys0}) and a Nambu-outer ordering \code{"tsltsl"}. The matrices
\code{tysx} $=\tau_y\sigma_x$ and \code{tys0} $=\tau_y\sigma_0$ are defined in cell 2, so every branch runs; only the default is used here.
\end{function}

\begin{function}{takagi(A)}
\begin{Verbatim}
def takagi(A):
    """D, U = takagi(A)
    Diagonalize a complex symmetric matrix as A = U D U^T
    where U is unitary and D is diagonal."""
\end{Verbatim}
Takagi factorization of a complex symmetric matrix. \emph{Method (from the code and its comments):} (i) $B = A^\dagger A$ is Hermitian
positive semidefinite; diagonalize $B = V b V^\dagger$. (ii) $C = V^TAV$ is complex symmetric with $C = X + iY$, $[X, Y] = 0$; diagonalize
$X + Y$ (\code{eigh}) to get $W$, a simultaneous eigenbasis, and set $U = VW$. (iii) $z_i = (U^TAU)_{ii} = r_ie^{i\theta_i}$; multiply column $i$ of
$U$ by $e^{-i\theta_i/2}$ (\code{(r/z)**0.5}) so the diagonal becomes real and non-negative. Returns \code{r} and \code{U.conj()}.

\emph{Why it works} (now in the code comments). $C^\dagger C=V^\dagger A^\dagger AV=\mathrm{diag}(b)$ is real, and with $C=X+iY$ ($X$, $Y$ real
symmetric) $C^\dagger C=X^2+Y^2+i[X,Y]$, so $[X,Y]=0$; one real orthogonal $W$ (the eigenvectors of $X+Y$, barring accidental degeneracy)
diagonalizes both, and $U^TAU=W^TCW$ is diagonal. After the phase fix $U^TAU=\mathrm{diag}(r)$, i.e.\ $A=U^*\mathrm{diag}(r)U^\dagger$: the
factor $Z$ of $A=ZDZ^T$ is $U^*$, hence \code{U.conj()}.
\end{function}

\begin{function}{individual\_MZMs(E, M, tol=1e-3, ordering, spinor)}
\begin{Verbatim}
    mask = np.abs(E) < tol
    psi = M[:, mask]
    C = dag(psi) @ PHconj(psi, ordering=ordering, spinor=spinor)
    c, Z = takagi(C)
    gamma = psi @ Z
    x = np.repeat(xs, 4)   # x coordinate of each component (site-major ordering, see the markdown)
    R = dag(gamma) * x @ gamma
    r, O = np.linalg.eigh(R.real)  # Want O ortho to conserve C
    mbs = gamma @ O
\end{Verbatim}
Separates individual Majoranas in the zero-energy subspace. (1) The subspace is $|E| < $ \code{tol} $=10^{-3}$; if it has an odd number of
states the closest state outside is added with a warning. (2) $C_{ij} = \langle\psi_i|C\psi_j\rangle$ is the (complex symmetric) matrix of the
antiunitary $C$ in the subspace; its Takagi vectors $Z$ give $\gamma = \psi Z$, eigenstates of $C$ (Majoranas): with $C_{\rm mat}=ZDZ^T$,
$C\gamma=(C\psi)Z^*=\psi C_{\rm mat}Z^*=\psi ZD(Z^TZ^*)=\gamma D$, and $D=\mathbb 1$ for a subspace closed under $C$. (3) The position operator
$R = \gamma^\dagger x\gamma$ is diagonalized with its real part ($R$ is real for Majoranas because $C$ commutes with $x$), so $O$ is real orthogonal and the $\gamma O$ stay eigenstates of $C$; the
Majoranas come out localized and sorted by $x$. Returns \code{mbs, mask, M}.

\emph{Change for this geometry} (from the markdown): the original localizes with \code{x = np.tile(range(N//4), 4)}, the site index. For the
site-major ordering \code{4*site + c} this should be \code{np.repeat}, and here the site index is not ordered along the chain, so the operator
uses the $x$ coordinate of each site, \code{np.repeat(xs, 4)}. The \code{ordering} and \code{spinor} arguments are passed on to \code{PHconj}; the
default spinor is \code{"udd-u"}, kwant's, as in the call (the copied code had \code{"uddu"} as default and did not pass the arguments on). Output: 4 states, centres
$x = -46.6, -17.5, 16.5, 45.6$.
\end{function}

\begin{function}{fix\_unitary(mbs, mask, M)}
\begin{Verbatim}
    U = np.array([[1, 1], [1j, -1j]]) / 2**0.5
    phi = mbs @ np.kron(np.eye(n // 2), U)
    ind = np.arange(n).reshape(n // 2, 2).swapaxes(0, 1).ravel()
    phi = phi[:, ind]
    ...
    ind = np.arange(n).reshape(2, n // 2); ind[0] = np.flip(ind[0]); phi = phi[:, ind.ravel()]
    M[:, mask] = phi
\end{Verbatim}
Writes fermions into the zero-energy columns so that the full unitary has the BdG form $M = \begin{psmallmatrix}U & V^*\\ V & U^*\end{psmallmatrix}$.
Each consecutive Majorana pair $(\gamma_{2j-1},\gamma_{2j})$ becomes $(\gamma_{2j-1}\pm i\gamma_{2j})/\sqrt2$ (the columns of
$U_2 = \tfrac1{\sqrt2}\begin{psmallmatrix}1&1\\ i&-i\end{psmallmatrix}$). The first reordering (\code{[0,2,1,3]} for four Majoranas) groups
the ``$+$'' combinations before the ``$-$'' combinations; the second reverses the ``$+$'' group. For four Majoranas the zero-energy columns
$2N-2,\dots,2N+1$ (the four middle eigenvectors of \code{eigh}) become
\[
\bigl(\tfrac{\gamma_3+i\gamma_4}{\sqrt2},\ \tfrac{\gamma_1+i\gamma_2}{\sqrt2},\ \tfrac{\gamma_1-i\gamma_2}{\sqrt2},\ \tfrac{\gamma_3-i\gamma_4}{\sqrt2}\bigr),
\]
so column $2N+k$ and column $2N-1-k$ are particle--hole partners ($C\gamma = \gamma$). The notebook calls columns $2N$, $2N+1$ $\psi_A$, $\psi_B$
and requires them to have positive energy. With the Python Majoranas, $\psi_A$ is therefore $(\gamma_1 - i\gamma_2)/\sqrt2$; the TN code writes
$(\gamma_1+i\gamma_2)/\sqrt2$ with its own sign of $\gamma_2$, which is a sign convention for $\gamma_2$ only (both are the positive-energy
fermion of pair A). The reason for the second reordering (now in the code comment): in the \code{eigh} ordering of a BdG matrix,
$E_{2N+k}=-E_{2N-1-k}$ and these two columns are $C$ partners, so placing $\varphi_k$ in column $2N-1-k$ and $C\varphi_k$ in column $2N+k$ keeps
the fixed $M$ in BdG form and makes the masks (column $2N+k$ occupied $\Leftrightarrow$ column $2N-1-k$ empty) valid for any number of pairs.
(The copied code computed the Majorana-to-fermion block twice; the duplicate has been removed.) \code{mbs2 = mbs[:, [0,1,2,3]]}
pairs $(\gamma_1,\gamma_2)$ and $(\gamma_3,\gamma_4)$. Checks printed: column energies $(-1.158, -1.090, 1.090, 1.158)\times10^{-9}$ (an
\code{assert} requires the last two to be positive; otherwise the order within a pair must be swapped), $\|U^\dagger U - 1\| = 6.5\times10^{-14}$,
$\max\|C\psi_{2N+k} - \psi_{2N-1-k}\| = 2.2\times10^{-14}$. \code{psi0 = U\_fixed.copy()} is the initial full basis.
\end{function}

\subsubsection*{Logical states}

\code{masks} is a $4\times4N$ boolean array, one row per state $\ket{n_An_B}$: $\ket{00}$, $\ket{01}$, $\ket{10}$, $\ket{11}$. Columns $2N, 2N+1$
are set to the occupations of $\psi_A$, $\psi_B$; all higher columns are empty; the lower half is set by
\code{masks[:, :2N] = np.flip(\textasciitilde masks[:, 2N:], axis=1)}, so a column below $2N$ is occupied exactly when its particle--hole partner
is empty. Thus $\ket{00}$ occupies all $2N$ negative-energy columns (including $C\psi_A$, $C\psi_B$), and $\ket{11}$ occupies $\psi_A,\psi_B$
instead of $C\psi_A, C\psi_B$. The logical qubit is $\ket0_L = \ket{00}$ (row 0), $\ket1_L = \ket{11}$ (row 3). \code{Z0 = psi0[:, 2N-2:2N+2]} is
the initial zero-mode subspace.

\subsubsection*{Sparse $H(t)$ and the Gauss--Legendre step}

\begin{function}{exchange\_sparse(f, J), ham\_sp(f)}
Sparse (CSC) versions: \code{exchange\_sparse} is a \code{block\_diag} of the $4\times4$ exchange blocks; \code{ham\_sp(f) = H\_noJ\_sp + exchange\_sparse(f)}.
\end{function}

\begin{function}{step(psi, t, dt, ft)}
\begin{Verbatim}
def step(psi, t, dt, ft):
    M0 = -1j * ham_sp(ft(t + c1*dt))
    M1 = -1j * ham_sp(ft(t + c2*dt))
    A = sp.bmat([[I_sp - 0.25*dt*M0, -a12*dt*M0],
                 [-a21*dt*M1, I_sp - 0.25*dt*M1]], format="csc")
    B = np.vstack([M0 @ psi, M1 @ psi])
    lu = splu(A, permc_spec="MMD_AT_PLUS_A")
    ...
    return psi + 0.5*dt*(K[0] + K[1])
\end{Verbatim}
One step of the 2-stage Gauss--Legendre implicit Runge--Kutta method (4th order) for $\dot\psi = M(t)\psi$, $M = -iH(t)$, applied to all
$4N$ columns at once. With nodes $c_{1,2} = \tfrac12\mp\tfrac{\sqrt3}{6}$ and Butcher coefficients $a_{11} = a_{22} = \tfrac14$,
$a_{12} = \tfrac14-\tfrac{\sqrt3}{6}$, $a_{21} = \tfrac14+\tfrac{\sqrt3}{6}$, $b_1 = b_2 = \tfrac12$, and $M_i = M(t + c_i\,dt)$ ($H$ at the two Gauss
points), the stages solve the linear system
\[
\begin{pmatrix} 1 - a_{11}\,dt\,M_0 & -a_{12}\,dt\,M_0\\ -a_{21}\,dt\,M_1 & 1 - a_{22}\,dt\,M_1\end{pmatrix}
\begin{pmatrix}K_0\\K_1\end{pmatrix} = \begin{pmatrix}M_0\psi\\ M_1\psi\end{pmatrix},
\qquad \psi(t+dt) = \psi + \tfrac{dt}{2}(K_0 + K_1).
\]
The $2D\times2D$ sparse matrix is factorized once per step by \code{splu} (permutation \code{MMD\_AT\_PLUS\_A}), and the $D$ right-hand sides are
solved in \code{n\_threads} chunks on a thread pool; \code{n\_threads = 1}, so the timings of section 11 are single-thread (the comment notes
that more threads are faster). This is the sparse-LU form of the period-finding notebook. Time step \code{DT = 0.25}.
\end{function}

\subsubsection*{Readout}

\code{hz} indexes the horizontal chain ($y = 0$) sorted by $x$; \code{tauz} $= (1,1,-1,-1)$ repeated over sites.

\begin{function}{electron\_density(W)}
\begin{Verbatim}
    """Electron number per site of the state with occupied columns W: electron part of diag(G), G = W W†."""
    return (np.abs(W) ** 2).sum(axis=1).reshape(N, 4)[:, :2].sum(axis=1)
\end{Verbatim}
$n(\mathbf r) = \sum_{\sigma}G_{\mathbf r\sigma,\mathbf r\sigma}$ with $G = WW^\dagger$ restricted to the electron components $(u_\uparrow, u_\downarrow)$.
\end{function}

\begin{function}{readout(psi)}
For the current full basis \code{psi} returns a dict with
\code{dn\_h}: $\delta n = n(\text{\code{psi}}[:, \text{\code{masks[3]}}]) - n(\text{\code{psi}}[:, \text{\code{masks[0]}}])$ on the horizontal chain;
\code{dQ}: $\sum_{\mathbf r}\delta n$ over all sites; \code{dQ\_modes}: $\langle\psi_k|\tau_z|\psi_k\rangle$ for columns $2N$, $2N+1$;
\code{z}: the $4\times2$ overlaps $Z_0^\dagger\psi_{A,B}$ with the initial zero-mode subspace. Because the two masks differ only in
columns $2N-2\dots2N+1$, \code{dQ} equals the sum of \code{dQ\_modes} (Overview, Logical qubit and readout); the bulk columns cancel exactly.
\end{function}

\begin{function}{evolve(ft, psi, t0, t1, dt\_snap, dt=DT)}
Evolves the full basis from $t_0$ to $t_1$ with \code{step}, renormalizing every column after each step, and calls \code{readout} at $t_0$, every
\code{dt\_snap} and at the last step. Returns \code{psi, ts, R} (list of readout dicts).
\end{function}

\begin{function}{mode\_leak(R, name)}
Leakage at the last snapshot: $\ell_k = 1 - \|z_{:,k}\|^2$ from \code{R[-1]["z"]}; prints and returns it. Since \code{Z0} is an orthonormal basis of
the same subspace as the four $\gamma_i(0)$, this equals Julia's \code{mode\_leak(z)}. Values of order $-10^{-15}$ are rounding.
\end{function}

\begin{function}{dQ\_adiabatic(H)}
Diagonalizes $H$ and sums $\langle\phi|\tau_z|\phi\rangle$ over the positive-energy states among the four lowest $|E|$, skipping states with
$|E| < 10^{-11}$ (an exact zero mode contributes 0). Evaluated at every fusion snapshot to draw the adiabatic curve in the cell-19 plot.
\end{function}

\paragraph{End of the fusion.} \code{q = Fusion()}, \code{Tf = 380}; \code{HT\_end = ham(f\_fuse(q, Tf))}; \code{Z4\_end} are the four eigenvectors of
smallest $|E|$. \code{TIMES} collects wall times.

\begin{function}{run(ft, T1, label, dt\_gate=5.0, dt\_fuse=1.0)}
One full run starting from \code{psi0}: if \code{T1 > 0}, two \code{evolve} calls of duration $T_1$ (snapshots every 5) with \code{mode\_leak}
after each (\code{leak\_one}, \code{leak}); then the fusion from $T_g = 2T_1$ to $T_g + T_f$ (snapshots every 1). Fusion leakage
$1 - \|Z_{4,\rm end}^\dagger\psi_k\|^2$ (the counterpart of Julia's \code{fusion\_leak}). Returns \code{out} with \code{Tg}, the leakages,
\code{out["fusion"]} (final \code{dQ}, \code{dQ\_modes}, \code{leak}, \code{ts}, \code{dQ\_t}, \code{dQad\_t}) and \code{out["snap"]} = (times,
$\delta n$ on the horizontal chain, $\Delta Q$) at all snapshots. Records the gate and fusion wall times in \code{TIMES}.
\end{function}

%----------------------------------------------------------------------
\subsection{Sections 7--8: Z and X gates, then the partial fusion}

The texture functions are \code{fz} (two \code{exchange\_12}, then \code{f\_fuse}), \code{ff} (fusion only) and \code{fx} (two
\code{exchange\_23} with \code{rot\_T = 400}, then fusion). A check prints $\|H(2T_{12}) - H(0)\| = 0$ with $T_{12} = 2040$: the texture returns
to the initial one after the gate. Results:

\begin{center}
\footnotesize\setlength{\tabcolsep}{3pt}
\begin{tabular}{lll}
\toprule
Run & Leakage $(\psi_A, \psi_B)$ & Fusion: $\Delta Q$ ($\psi_A$, $\psi_B$); leakage\\
\midrule
Z, one exchange & $(9.74\times10^{-7}, 0)$ & \\
Z, two exchanges & $(1.997\times10^{-6}, 0)$ & $-0.477717$ ($-0.477717$, $3.8\times10^{-8}$); $(3.6963\times10^{-5}, 0)$\\
F (no gate) & & $-0.477373$ ($-0.477373$, $3.8\times10^{-8}$); $(3.7778\times10^{-5}, 0)$\\
X, one exchange & $(3.4886\times10^{-5}, 1.24\times10^{-7})$ & \\
X, two exchanges & $(3.4163, 3.4869)\times10^{-5}$ & $+0.477004$ ($+0.476113$, $8.9\times10^{-4}$); $(8.075, 3.331)\times10^{-5}$\\
\bottomrule
\end{tabular}
\end{center}

\begin{figure}[H]
\centering
\includegraphics[width=0.9\textwidth]{py_dQ.png}
\caption{Python, cell 19: $\Delta Q$ during the fusion. Left: after Z and with no gate, against the adiabatic value
(\code{dQ\_adiabatic}, dotted). Right: after X, against minus the adiabatic value.}
\end{figure}

%----------------------------------------------------------------------
\subsection{Section 9: charge density \texorpdfstring{$\delta n(x,t)$}{δn(x,t)}}

\begin{function}{density\_figure(ts, x, D, Tg, gate, method, out)}
The matplotlib counterpart of the Julia \code{density\_figure}: $\log_{10}|\delta n|$ over the full protocol (range $[-12, 0]$, dashed at
$T_g/2$ and $T_g$) and $\delta n$ during the fusion on a symmetric scale; dotted lines at the junctions. Saved to \code{out}.
\end{function}

The snapshot data of the Z and X runs are saved to \code{gate\_density\_kwant\_\{Z,X400\}.npz} (\code{ts, x, D, Tg, dQ}; \code{dQ} summed over
all sites). Z: $T_g = 4080$, 1197 snapshots, $\Delta Q(T) = -0.477717$; X: $T_g = 7760$, 1933 snapshots, $\Delta Q(T) = +0.477004$.

\begin{figure}[H]
\centering
\includegraphics[width=\textwidth]{py_density_X.png}
\caption{Python, X run (\code{density\_figure}). Left: $\log_{10}|\delta n(x,t)|$ over two \code{exchange\_23} ($T_1 = 3880$) and the fusion.
Right: $\delta n$ during the fusion, with the opposite sign to the Z run ($\sum_{\rm chain}\delta n = +0.444$).}
\end{figure}

%----------------------------------------------------------------------
\subsection{Section 10: saving the results}

Top-level code: \code{iA} is the $+E_A$ state among the four lowest $|E|$ of \code{HT\_end}; the dict \code{gates} holds the same 12 keys as
\code{msh\_tjunctions\_tn\_gates.csv}, with \code{fus\_dQad\_A} $=\langle\phi_{iA}|\tau_z|\phi_{iA}\rangle = -0.476854$ and
\code{fus\_E\_A} $= 0.0438539$. Saved to \code{gates\_kwant.npz}.

%----------------------------------------------------------------------
\subsection{Section 11: cost of the dense route}

Measured times (one thread, $dt = 0.25$): Z gate 570.6~s (16\,320 steps, 34.96~ms/step), Z fusion 54.4~s (1520 steps), F fusion 54.0~s,
X gate 1110.0~s (31\,040 steps, 35.76~ms/step), X fusion 54.1~s; one \code{eigh} of $H$ ($D = 576$) 0.078~s.

\begin{function}{ham\_nleg(nleg, f=f0)}
Builds the same array with arms of \code{nleg} sites (junctions at $(c-2)\cdot2\,$\code{nleg}) as a kwant system and returns its sparse BdG
matrix and the number of sites. Self-contained copy of the label, texture and kwant builder for arbitrary leg length.
\end{function}

\begin{function}{bench\_step(H, dt=DT, nst=3)}
Average time of one full-basis Gauss--Legendre step with a static $H$, as in \code{step}: builds $M_0, M_1$ (as \code{(H + 0*H).tocsc()}, so
a sum of the static $H$ with itself, not the rebuild of $H(t)$ by \code{ham\_sp}), the block matrix, the sparse LU and the solve for
$D$ right-hand sides (\code{psi = I}).
\end{function}

\begin{center}
\begin{tabular}{rrrrr}
\toprule
Arm & sites & $D$ & full-basis step & \code{eigh}\\
\midrule
16 & 144 & 576 & 24.6 ms & 0.078 s\\
32 & 288 & 1152 & 157.6 ms & 0.544 s\\
64 & 576 & 2304 & 1095.3 ms & 5.649 s\\
\bottomrule
\end{tabular}
\end{center}
Everything is saved to \code{timing\_kwant.npz} (\code{time\_*}, \code{steps\_*}, \code{dt}, \code{eigh}, \code{bench\_*}). The benchmark step at
144 sites (24.6~ms) is faster than the measured step in the runs (about 35~ms). The markdown now explains why: a step of \code{evolve} also
rebuilds the sparse $H(t)$ at both Gauss points (\code{ham\_sp}), goes through the thread-pool solve (which copies $B$ into contiguous chunks
and stacks the result, even with one thread) and normalizes $\psi$, plus a readout every few steps. Timed one at a time on the same machine
(under some load, so all times were longer), they cost about 5, 7 and 1~ms per step, which accounts for the gap. The benchmark measures how
the step arithmetic scales with $D$; the extra work is overhead of this implementation, not of the method.
